\documentclass[DIN, pagenumber=false, fontsize=11pt, parskip=half]{scrartcl}
\usepackage[UTF8]{ctex}  % 中文支持
% \usepackage{ngerman}
\usepackage[utf8]{inputenc}
\usepackage[T1]{fontenc}
\usepackage{textcomp}
\usepackage{graphicx}
\usepackage{amsmath}
\usepackage{booktabs}
\usepackage{tasks}
\usepackage{xcolor} % Required for custom red
\usepackage{amsmath,amssymb}
\usepackage{tikz}
\usetikzlibrary{calc}
\usepackage{geometry} % Custom margins

\usepackage{tocloft} % Modifying the table of contents

% for matlab code
% bw = blackwhite - optimized for print, otherwise source is colored
\usepackage[framed,numbered,bw]{mcode}
\usepackage[most]{tcolorbox}
% for other code
%\usepackage{listings}

\setlength{\parindent}{0em}

% set section in CM
\setkomafont{section}{\normalfont\bfseries\Large}

\newcommand{\mytitle}[1]{{\noindent\Large\textbf{#1}}}
\newcommand{\mysection}[1]{\textbf{\section*{#1}}}
\newcommand{\mysubsection}[2]{\romannumeral #1) #2}

\newtcolorbox{parchmentquote}{
    enhanced,
    colback=orange!5!white,   % 浅米黄色背景
    colframe=orange!50!black,
    boxrule=0.5pt,
    arc=2pt,                  % 细微圆角
    boxsep=5pt,
    before skip=10pt,
    after skip=10pt,
    fontupper=\itshape,
    borderline west={2pt}{0pt}{orange!80!black} % 加厚的左装饰线
}
%===================================
\begin{document}
\newcommand{\blank}[1][3em]{\underline{\hspace{#1}}}
\noindent\textbf{课后练习} \hfill \textbf{online course}\\
primary school \hfill Liu\\

\mytitle{分数练习 \hfill \today}
\newline
\begin{parchmentquote}
    比例是分数的综合应用，难度比较大


\end{parchmentquote}
\begin{enumerate}
  \item 下列各比，能与 \(6 : 5\) 组成比例的是（  \quad   ）。

A. \(\frac{1}{3} : \frac{2}{5}\) \quad
B. \(\frac{1}{5} : \frac{1}{6}\) \quad
C. \(\frac{1}{6} : \frac{1}{5}\) \quad
D. \(\frac{2}{5} : \frac{3}{2}\)
\item 一款手表的零件长 \(0.5\) 毫米，在设计图纸上的长度是 \(10\) 厘米，图纸的比例尺是（  \quad   ）。

A. \(1 : 20\) \quad
B. \(20 : 1\) \quad
C. \(1 : 200\) \quad
D. \(200 : 1\)
\item 照片上悦悦的身高是 \(4\) cm，她的实际身高是 \(1.6\) m。这张照片的比例尺是（  \quad   ）。

A. \(1 : 0.4\) \quad
B. \(1 : 40\) \quad
C. \(0.4 : 1\) \quad
D. \(40 : 1\)
\item 一个底边是 \(8\) 厘米、高是 \(4\) 厘米的平行四边形，按 \(2 : 1\) 的比放大后的平行四边形面积与原来的平行四边形面积之比是（  \quad   ）。

A. \(1 : 2\) \quad
B. \(2 : 1\) \quad
C. \(1 : 4\) \quad
D. \(4 : 1\)
\item 下面不能用比例 “\(3 : 6 = 9 : x\)” 解决的是（  \quad   ）。

A. 将一根圆木锯成 \(3\) 段需要 \(6\) 分钟，照这样算，锯成 \(9\) 段需要 \(x\) 分钟 \\
B. 一辆汽车 \(3\) 分钟行驶 \(6\) 千米，照这样算，\(9\) 分钟可以行驶 \(x\) 千米 \\
C. \(3\) 支铅笔可以换 \(6\) 块橡皮，照这样算，\(9\) 支铅笔可以换 \(x\) 块橡皮 \\
D. 长 \(6\) 厘米、宽 \(3\) 厘米的长方形按一定的比画出新图后，长是 \(x\) 厘米，宽是 \(9\) 厘米
\item 下面每个选项中的两个比不能组成比例的是（  \quad   ）。

A. \(4 : 3\) 和 \(60 : 45\) \quad
B. \(1 : \frac{1}{5}\) 和 \(\frac{5}{8} : \frac{1}{8}\) \quad
C. \(1 : 1.2\) 和 \(3.5 : 4.2\) \quad
D. \(20 : 10\) 和 \(0.5 : 1\) 
\item\[
\begin{array}{|c|c|c|c|c|}
\hline
\text{甜度} & \text{无糖} & \text{三分糖} & \text{五分糖} & \text{七分糖} \\
\hline
\text{糖 / 克} & 0 & 2.4 & 4 & 5.6 \\
\hline
\text{奶茶 / 毫升} & 100 & 100 & 100 & 100 \\
\hline
\end{array}
\]

(1) 悦悦准备按上面的配比表配制一杯中杯（500毫升）的五分糖奶茶，她需要准备多少克糖？（用方程解）

\vspace{1cm}

(2) 《中国居民膳食指南（2022）》中指出“每个人每天糖的摄入量不要超过50克”。按照这种膳食指导，一个人每天最多能喝多少毫升三分糖奶茶或多少毫升七分糖奶茶？（得数保留整百数）
\end{enumerate}








\end{document}